\section{From Total Rankings to Certified Partial Orders}
\label{sec:ranking}

\subsection{Prequential forecasts and score bands}

At month $t$, let $\cA_t=\{1,\ldots,N_t\}$ be the eligible stock universe. A predetermined learning algorithm produces the next-month return forecast $\widehat\mu_{i,t}$ using information available at $t$. The algorithm may be linear, tree based, or a neural network. The primary implementation uses a fixed-length rolling training window, so the forecast sequence is genuinely out of sample and has finite memory.

Let $\widehat s_{i,t}>0$ be a reliability scale. It can be learned from prequential absolute residuals, model disagreement, or both, but it is computed without using $R_{i,t+1}$. The paper does not interpret
\[
[\widehat\mu_{i,t}-q\widehat s_{i,t},
  \widehat\mu_{i,t}+q\widehat s_{i,t}]
\]
as a stock-specific confidence interval. The bands are a nested device for deciding which relative comparisons the investor is willing to assert. Their calibration concerns the realized ranking loss of the complete decision rule.

For $q\geq0$, define
\begin{equation}
\underline\mu_{i,t}(q)
=
\widehat\mu_{i,t}-q\widehat s_{i,t},
\qquad
\overline\mu_{i,t}(q)
=
\widehat\mu_{i,t}+q\widehat s_{i,t}.
\label{eq:bands}
\end{equation}

\begin{definition}[Certified rank relation]
For distinct stocks $i,j\in\cA_t$, write
\begin{equation}
i\succ_{q,t}j
\quad\Longleftrightarrow\quad
\underline\mu_{i,t}(q)
>
\overline\mu_{j,t}(q).
\label{eq:relation}
\end{equation}
When neither $i\succ_{q,t}j$ nor $j\succ_{q,t}i$ holds, the rule abstains on the pair.
\end{definition}

Let $A_{ij,t}(q)=\ind\{i\succ_{q,t}j\}$. The next proposition gives the basic geometry.

\begin{proposition}[Nested strict partial rankings]
\label{prop:partial-order}
For every $t$ and $q\geq0$, the relation $\succ_{q,t}$ is irreflexive, asymmetric, and transitive. Moreover, if $q'\geq q$, then
\[
A_{ij,t}(q')
\leq
A_{ij,t}(q)
\]
for every pair $(i,j)$.
\end{proposition}

The result means that abstention never reverses a rank relation. It only deletes relations that are not sufficiently separated. Transitivity follows from interval separation, which is why we use stock-level bands instead of an arbitrary pairwise margin that could generate cycles.

\subsection{Reliable breadth}

Let $\omega_{ij,t}=\omega_{ji,t}\geq0$ be a pair weight measurable at $t$. The baseline uses equal weights; alternatives emphasize top-versus-bottom relations, market capitalization, or trade capacity. Define the total weight of the raw ranking by
\[
W_t
=
\sum_{i\neq j}
\omega_{ij,t}A_{ij,t}(0).
\]
Absent forecast ties, exactly one direction is active for each unordered pair with positive weight.

\begin{definition}[Reliable breadth]
The reliable breadth of the cross-section at threshold $q$ is
\begin{equation}
\cB_t(q)
=
\frac{
\sum_{i\neq j}\omega_{ij,t}A_{ij,t}(q)
}{W_t},
\label{eq:reliable-breadth}
\end{equation}
with $\cB_t(q)=0$ when $W_t=0$.
\end{definition}

Reliable breadth lies in $[0,1]$, equals one at $q=0$ under no ties, and is nonincreasing in $q$. It measures relations, not securities. A cross-section with one thousand stocks can have high nominal breadth and low reliable breadth if most forecast differences are small relative to their local scale.

The empirical analysis reports several versions. All-pair breadth measures the entire ordering. Tail breadth restricts $\omega_{ij,t}$ to pairs in which one stock lies in the raw top quantile and the other in the raw bottom quantile. Rank-distance breadth groups pairs by the distance between their raw ranks. These distinctions determine whether unreliability is confined to the middle of the cross-section or reaches the positions that drive conventional long--short strategies.

\subsection{Pairwise decision losses}

The realized error of an active relation is
\[
M_{ij,t+1}
=
\ind\{R_{i,t+1}\leq R_{j,t+1}\}.
\]
For candidate $q$, define the selected-pair error rate
\begin{equation}
e_t(q)
=
\frac{
\sum_{i\neq j}\omega_{ij,t}A_{ij,t}(q)M_{ij,t+1}
}{
\sum_{i\neq j}\omega_{ij,t}A_{ij,t}(q)
},
\label{eq:selected-error}
\end{equation}
with $e_t(q)=0$ when no relation is active. This is the natural answer to: among the comparisons the investor was willing to use, what fraction were wrong?

The denominator in \eqref{eq:selected-error} changes with $q$, so $e_t(q)$ need not be monotone. Directly applying an ordered risk-controlling selector would therefore be invalid. For an increasing finite candidate grid $q_1<\cdots<q_K$, define
\begin{equation}
L_t(q_k)
=
\max_{\ell\geq k}e_t(q_\ell).
\label{eq:monotone-envelope}
\end{equation}
Then $0\leq e_t(q_k)\leq L_t(q_k)\leq1$ and $L_t(q_k)$ is nonincreasing in $k$. We certify the stationary risk
\begin{equation}
\cR(q_k)
=
\E[L_t(q_k)].
\label{eq:ranking-risk}
\end{equation}
Controlling \eqref{eq:ranking-risk} controls the expected selected-pair error at the deployed threshold and protects against accidental nonmonotonicity at stricter thresholds.

A secondary loss avoids the random denominator:
\begin{equation}
L_t^{\mathrm{mass}}(q)
=
\frac{
\sum_{i\neq j}\omega_{ij,t}A_{ij,t}(q)M_{ij,t+1}
}{W_t}.
\label{eq:wrong-mass}
\end{equation}
It is exactly nonincreasing in $q$ and measures the weighted mass of wrong bets per raw opportunity. The main text emphasizes \eqref{eq:selected-error} because its interpretation is closer to ranking reliability; \eqref{eq:wrong-mass} is a robustness target.
